%% 第四章

\chapter{自回归生成模型与 B-rep 重建方法}
\chaptermark{自回归生成模型与 B-rep 重建方法}

本章描述生成模型的总体架构、损失与优化、推理采样以及从 token 序列到 STEP / STL 文件的逆向重建管线，整体顺序由模型训练自然过渡到生成与重建评测。主实验采用“窗口嵌入层 + LLaMA3 风格因果主干 + 线性输出头”的轻量自回归结构；代码中另保留未启用的窗口化解码扩展。下文重点说明主实验所用结构，并在 §4.4 单独说明扩展模块边界，以便后续工作能在不混淆主分支的前提下接入。

\begin{figure}[!htbp]
    \centering
    \includegraphics[width=0.82\textwidth]{插图/fig4-1_model_architecture.pdf}
    \caption{自回归生成模型总体架构示意图。输入 token 序列首先经过 token embedding 与 RoPE 位置编码，随后进入由多层因果 Transformer block 组成的 LLaMA3 风格主干，每个 block 包含 RMSNorm、Grouped Query Attention 与 SwiGLU 前馈层。主干末端经最终 RMSNorm 和线性输出层得到下一 token 的 logits，并据此采样得到下一个 token。}
    \label{fig:model-architecture}
\end{figure}

\section{模型总体架构}

\subsection{窗口化输入与生成分支}

模型把第~3 章生成的分组序列张量（形状 $(L, 3)$）按窗口大小 $W$ 进一步组织为 $(L', W)$ 形式，其中 $L'$ 是窗口数。窗口大小控制模型一次预测所覆盖的 token 数：当 $W = 1$ 时，模型退化为“标准 next-token”自回归，每一行作为单步预测目标，主干 Transformer 的最终隐藏状态直接经过线性输出层得到下一个 token 的 logits；当 $W > 1$ 时，可启用窗口内部解码器，对窗口内的 $W$ 个分量做第二层因果自回归。

本文主实验使用 $W = 1$ 配置，对应标准 next-token 输出头，不启用窗口内部解码器。该扩展分支仅作为后续工作保留，相关边界见 §4.4。

采用 $W = 1$ 的直接好处是模型行为容易验证。输入序列中的每个有效 token 都对应一个明确的下一 token 预测目标，训练损失可以按 token 平均直接解释为 token 级交叉熵，推理时每一步只新增一个 token，生成过程与第~3 章定义的 1D 序列完全一致，可以逐 token 与训练数据上的统计模式对齐。对于本文这种以工程闭环为主的本科毕设，先采用最朴素、最可复查的分支更有利于失败归因——如果生成结果失败，可以从 token 类型错误、END 终止异常、反量化坐标异常、NURBS 拟合失败或拓扑缝合失败等环节按顺序定位，而不必先把“窗口化分解是否引入了未知误差”作为变量排除。一旦标准 next-token 分支已经验证可行，再启用 $W > 1$ 时的提升或退化都能在同一基线上做对照。


\subsection{模型级联结构}

主实验中模型由两段级联构成：

\begin{enumerate}
    \item 窗口嵌入层：把窗口内 $W$ 个 token id 融合为一个 $d_{\text{model}}$ 维向量，
    \item LLaMA3 风格主干：$N$ 层因果 Transformer，对窗口向量序列做长程依赖建模，其末层经 RMSNorm 后接线性输出层，直接输出下一个 token 的 logits。
\end{enumerate}

从数据流上看，输入 token 首先被映射到连续向量空间，由因果 Transformer 在时间维上聚合历史上下文，最后再投影回模型词表。这一流程沿用了语言模型生成文字的基本范式，但本文 token 的语义比文本更混合：数据侧统一词表中既有量化坐标，又有面索引和结构控制符（START / SEP / END），三类 token 的数值范围、出现位置和上下游依赖都不同；模型侧词表则额外包含 PAD 这一填充 id。这意味着模型不仅要学习坐标分布，还要隐式地学习“在当前位置应该输出哪一类 token”以及“START / SEP / END 何时出现”。第~3 章中稳定的面段、边段骨架为这种学习提供了强先验：每隔 49 或 14 个 token 就会自然落到面索引 / 邻接面索引位置，每段之间用 SEP 切分，模型可以把“距离上一次 SEP 多少 token”作为辅助上下文，从而自发形成位置类型预测的近似规则。第~5 章中 hr2\_patch1 的 S1 直接清零，正是这一隐式语法学习在生成阶段的体现。

主干隐藏维度 $d_{\text{model}}=256$ 相比大语言模型明显更小，这是由数据规模、序列长度和硬件预算共同决定的——CAD token 序列虽然具有长程依赖，但训练样本规模远小于互联网文本，过大的模型会在有限数据上过拟合，也会显著增加 4\,096 长度下的显存消耗。本文采用 8 层、256 维的轻量配置，目标是让模型足以捕获 DeepCAD 子集中的主要结构规律，同时保证 4 张 H20 上可以完成 200 epoch 训练。

\section{窗口向量化嵌入}

\subsection{Token 嵌入与线性融合}

窗口嵌入层包含一张大小为 $V \times d_{\text{model}}$ 的 token Embedding 表 $E_{\text{val}}$ 与一层线性融合 $W_{\text{fuse}} \in \mathbb{R}^{(W \cdot d_{\text{model}}) \times d_{\text{model}}}$，其中 $V$ 为模型词表规模。给定窗口 $x_t \in \mathbb{Z}^{W}$，先按位查表得到 $W$ 个 $d_{\text{model}}$ 维向量并展平，再通过线性层投影回 $d_{\text{model}}$：

\begin{equation}
    \mathbf{h}_t = W_{\text{fuse}} \cdot \operatorname{concat}\big(E_{\text{val}}(x_t^{(0)}), E_{\text{val}}(x_t^{(1)}), \ldots, E_{\text{val}}(x_t^{(W-1)})\big).
    \label{eq:embed-fuse}
\end{equation}

在主实验 $W = 1$ 的设定下，融合层退化为单 token 的线性变换。当 $W = 3$ 时该结构可同时将一个三元组的三个分量打包成一个窗口向量，避免使用三张独立的 Embedding 表。

即使在 $W = 1$ 时，窗口嵌入层仍然承担两类工程作用。第一类是把离散 token id 映射到连续空间，让坐标 token、面索引 token 与结构 token 在同一 embedding 表中共享参数化向量，使 Transformer 后续可以用统一的注意力机制处理三类语义。共享同一张 embedding 表也意味着这三类 token 在 logit 头上是在同一概率单纯形上竞争，这就要求模型在每个位置上隐式判别该位置应该属于哪一类 token，而判别信号来自上下文模式。第二类作用是保留与窗口化分支兼容的入口接口：未来若启用 $W = 3$ 或更大的窗口，把一个三元组打包为单个窗口向量时，只需把已有的查表向量拼接后经同一份 $W_{\text{fuse}}$ 投影，主干结构完全不必改动。这种实现方式确实牺牲了一点结构极简性，但显著提高了代码复用性，也让本文实验与窗口化扩展之间能共享同一模型入口与同一份 checkpoint 加载逻辑。


\subsection{词表配置与填充符号}

主实验中模型词表规模为 4\,054，即数据词表 4\,053 加上 1 个 PAD。PAD 的 id 为 4053，保留在 Embedding 表中以便填充位置可参与查表；训练时这些位置的 label 设为 $-100$，由交叉熵损失忽略（详 §4.5）。

\section{LLaMA3 风格生成主干}

\subsection{Transformer 块结构}

主干由 $n_{\text{layers}} = 8$ 个同构 LLaMA3 风格 Transformer block 串联。每个 block 采用 pre-norm 残差结构：

\begin{align}
    \mathbf{h}' &= \mathbf{h} + \text{GQA}\big(\text{RMSNorm}(\mathbf{h})\big), \label{eq:block-attn} \\
    \mathbf{h}'' &= \mathbf{h}' + \text{SwiGLU}\big(\text{RMSNorm}(\mathbf{h}')\big). \label{eq:block-ffn}
\end{align}

注意力部分采用 Grouped Query Attention，前馈部分采用 SwiGLU，残差前后各使用一层 RMSNorm。本文主干配置如下：层数 $n_{\text{layers}} = 8$、查询头数 $n_{\text{heads}} = 8$、KV 头数 $n_{\text{kv\_heads}} = 4$（每个 KV 头由 2 个 Q 头共享）、隐藏维度 $d_{\text{model}} = 256$、前馈中间维度 $\text{dim\_ffn} = 1024$，RoPE 的 $\theta = 500\,000$，RMSNorm 的 $\epsilon = 10^{-5}$。

pre-norm 残差结构相对 post-norm 的主要优点是深层训练更稳定：每个子层在进入注意力或前馈网络之前先做归一化，残差分支保留原始隐藏状态，相当于把“残差路径”和“非线性变换路径”在数值尺度上分离开来，梯度可以沿残差路径直接回传到底层而不必穿过整层的非线性变换。本文主干虽然只有 8 层、并不属于极深网络，但序列长度较长（最长 4\,096）、token 类型混合，训练稳定性仍是不可忽视的工程指标——一次损失爆炸就可能导致几个 epoch 的算力浪费。RMSNorm 相较 LayerNorm 只保留均方根缩放、去掉了按 batch 的均值统计，不仅计算量略小，在 GPU 上的算子也更利于融合，同时与 LLaMA 系列结构保持一致，方便日后从已有 LLaMA 实现迁移或对比。本文训练日志全程未出现 NaN 或显著损失反弹，间接说明 pre-norm + RMSNorm 的组合在当前数据规模和混合精度配置下是稳健的。


\begin{figure}[!htbp]
    \centering
    \includegraphics[width=0.82\textwidth]{插图/fig4-2_gqa_vs_mha.pdf}
    \caption{多头注意力（MHA）与 Grouped Query Attention（GQA）的结构对比。MHA 为每个 Query 头配置独立的 Key / Value 头，GQA 将多个 Query 头划分为一组并共享同一组 Key / Value，从而在保持多查询表达能力的同时减少 KV cache 的存储与读取开销。}
    \label{fig:gqa-vs-mha}
\end{figure}

\subsection{RoPE 位置编码实现}

RoPE 仅对 Q、K 应用，不作用于 V。RoPE 缓存的长度由模型配置中的最大序列长度决定，并在遇到超长输入时自动扩展。RoPE 的相对位置不变性使得主干在窗口数较多的场景下仍可保持稳定的注意力点积。

\subsection{KV 缓存与训练加速机制}

主干在推理阶段使用 KV cache：每层缓存最近一次前向的 $K$、$V$，新增窗口只需与缓存做注意力点积，把单步推理复杂度由 $O(T^2)$ 降至 $O(T)$。训练阶段启用梯度检查点（gradient checkpointing）与 PyTorch 自动混合精度，在 H20 上以 FP16/BF16 算子加速矩阵乘累加，训练日志中记录了“梯度检查点已启用（节省显存）”。

\subsection{主干输出层设计}

主干末端为一层 RMSNorm，其输出的隐藏状态 $\mathbf{h}_t \in \mathbb{R}^{d_{\text{model}}}$ 直接送入一层无偏置线性输出层，映射到 4\,054 维 logits，预测下一个 token 的分布。该“主干归一化层 $\to$ 线性输出层”即本文主实验所采用的标准 next-token 头。

由于模型词表同时覆盖面索引、坐标、结构 token 以及 PAD，输出层的每一次预测都是在完整 4\,054 维词表上竞争；其中前 4\,053 个 id 对应可写入原始序列的数据 token，id 4053 对应模型侧填充。训练过程中，模型只看到 token id 与其上下文，没有任何显式的“位置类型”标签，它必须从训练数据的统计模式中自行学到：在 START 之后的 48 个 token 位置上几乎只能输出坐标段 token，第 49 个位置上几乎只能输出面索引段 token，达到一定长度后逐步出现 SEP，再以同样规律进入边段，直至 END 收尾。这种隐式语法学习让实现保持简洁，无需在解码时维护显式状态机或加掩码，但也意味着模型偶尔可能打破序列格式，例如在面段中突然吐出 SEP、在坐标位置上输出面索引等。第~5 章中 hr2\_patch1 在 1\,005 次采样中 S1 直接为 0，可以从侧面说明在当前数据规模和采样配置下，这种隐式语法的稳定性已经足以支撑可解析序列的批量生成。

\section{窗口化扩展模块边界}

\begin{sloppypar}
本文代码中保留了面向 $W > 1$ 窗口化分支的局部解码器与局部 Transformer 块。其原始设计意图是在一个窗口包含多个语义分量时，将主干隐藏状态作为上下文，再对窗口内部 token 做局部因果分解。但本文主实验配置为 $W = 1$，模型只构建标准 next-token 输出头，不构建也不调用窗口内部解码器。因此，本文报告的训练损失、生成有效率和重建结果均来自标准 next-token 分支，不依赖该扩展模块。
\end{sloppypar}

为避免将未启用模块误写为本文核心方法，本章不再单独给出窗口内部解码器结构图，也不对其效果作定量声明。后续若启用 $W > 1$ 的窗口化设定，可在相同数据划分与评测口径下补充对照实验，再讨论该模块对生成质量和有效率的影响，相关方向列入 §5.6 未来工作。

\section{损失函数与训练策略}

\subsection{交叉熵损失}

总损失为 $-100$ 屏蔽下的 token 级交叉熵：

\begin{equation}
    \mathcal{L} = -\sum_{t,\, \text{label} \ne -100} \log p_{\theta}(y_t \mid y_{<t}),
    \label{eq:ce-loss}
\end{equation}

训练目标采用交叉熵损失。PAD 位置的 label 在数据侧已置为 $-100$，并在损失计算中被跳过，因此不会产生梯度。

交叉熵损失是训练阶段最直接的优化目标，它给出“模型对真实下一 token 的对数概率”这一可优化标量，但它与最终 BRepCheck 通过率之间并非线性等价。较低的 token 预测损失通常意味着模型更接近训练数据分布，生成序列也更可能满足常见骨架——这一点在第~5 章 hr2\_patch1 验证损失 0.10486 与 of saved 97.51\% 同时出现的现象上得到了印证。但同一段损失下降中，不同位置 token 的预测误差对最终几何后果并不对称：坐标 token 的局部偏差只会带来轻微几何误差（在 4000 级量化下最多 $5 \times 10^{-4}$ 归一化误差，落到几何上不大可能破坏闭合），而 SEP、面索引、END 等关键结构 token 的错误会让整段序列要么提前截断、要么连接到错误的面，进而触发解析失败或拓扑不闭合。也就是说，CE 是按 token 平均的，而 BRepCheck 是按整体实体二值判定的。出于这一点，本文在第~5 章同时报告验证集 CE、PPL、生成有效率与三态分布，避免只用损失曲线推断工程可用性。

\subsection{优化器与学习率调度}

优化器为 AdamW\cite{kingma_2015_adam,loshchilov_2019_adamw}，初始学习率 $5 \times 10^{-4}$、权重衰减 $0.01$、dropout $0.1$，学习率调度采用“warmup $\to$ 中段稳定 $\to$ 末段衰减”三段式形态，对应曲线共有 201 条 step-level 记录。训练全程启用混合精度\cite{micikevicius_2018_mixed_precision}，并在主干上启用梯度检查点以节省显存。

AdamW 把权重衰减项从梯度更新中解耦出来，先按学习率更新参数，再单独乘以衰减因子，与 Adam 把权重衰减揉进 L2 项相比能更稳定地避免学习率—权重衰减之间的隐式耦合，因此被广泛用于 Transformer 类模型训练。本文的初始学习率 $5 \times 10^{-4}$ 相对常见大语言模型预训练（多取 $1 \sim 3 \times 10^{-4}$）并不算低，但考虑到本文模型规模较小（约 5M 参数级别）、数据集经过规则化处理且训练序列长度固定，较高的初始学习率仍能在 200 epoch 内带来稳定下降，未观察到任何损失反弹。dropout = 0.1 用于在有限数据上抑制过拟合，梯度检查点通过在反向传播时重算部分中间激活换取显存节省，是支撑 4\,096 长度训练在单卡 95 GB 上仍有充足显存余量的关键。混合精度训练降低了矩阵乘法的开销，但对几何生成任务而言仍需关注数值稳定性——本文训练日志中未观察到损失爆炸或 NaN，说明该配置在当前硬件与数据上是稳健的。

\subsection{分布式训练}

训练使用 PyTorch DDP，以 4 个进程对应 4 张 H20，每进程 batch size 为 32，全局 batch size 为 128。详细超参数、启动方式与硬件配置见 §5.1。

\subsection{训练监控与模型选择}

训练过程中同时记录 epoch 级验证集交叉熵、step 级学习率以及 best checkpoint 的更新情况。本文最终采用 hr2\_patch1 的 Epoch~200 best checkpoint，是因为该 checkpoint 在验证集上取得最低交叉熵，且训练曲线未出现后期反弹。选择 checkpoint 时没有使用生成有效率反向调参，避免把最终评测样本过早引入模型选择过程，生成与 BRepCheck 评测只在训练完成后执行。

从工程复现角度看，checkpoint 选择需要同时关注数值指标和文件一致性。数值上，验证集 CE 越低通常说明模型对 token 序列分布拟合越好；文件上，checkpoint 必须与训练时的模型配置、词表规模、PAD id 和窗口大小完全一致，否则加载后即使参数维度匹配，也可能因 token 语义不一致导致生成失败。本文在训练日志和第~5 章表格中重复列出这些关键配置，就是为了让读者能够追溯“某个有效率数字”对应的具体模型状态。

此外，本文没有把训练损失最低视为 checkpoint 选择的唯一目标。若某个 checkpoint 验证 CE 略低，但采样时频繁提前输出 END 或输出不完整边段，它未必能带来更高的 BRepCheck 通过率——这正是 CE 与 BRepCheck 之间非线性关系的另一种体现。理想情况下，应在固定验证集损失之外，额外维护一组小规模生成验证样本，周期性统计 S1 / S2 / S3 变化并写入 TensorBoard，把“采样—重建—校验”作为训练过程中的轻量指标。本文受训练周期与显卡时间限制，仅在最终 checkpoint 上做一次完整生成评测，没有把生成—校验循环写进训练 loop，因此该自动化监控被列入后续改进方向，由 §5.6 中“失败案例分类”一并承担。

\section{自回归推理生成方法}

\subsection{采样策略}

推理时输入只含一个 START 三元组，模型按因果顺序逐窗口生成，每一步从主干末端的线性输出层取得 logits 后做温度缩放，再依次按 top-k 与 top-p 过滤，最后从过滤后的分布中随机采样。本文实际使用的采样配置为 top-k = 1 + top-p = 0.9（贪心式采样并保留累积概率 90\% 的候选集合），该类截断采样思路与神经文本生成中的 nucleus sampling 口径相关\cite{holtzman_2020_nucleus}。

伪代码如\autoref{alg:ar-generate} 所示。

\begin{algorithm}[htbp]
    \caption{自回归序列生成（$W = 1$ 主干 next-token 分支）}
    \label{alg:ar-generate}
    \begin{algorithmic}[1]
        \Require 训练好的模型 $M$，最大窗口数 $L_{\max}$
        \Ensure token 序列 $x$
        \State $x \gets [\text{START}]$，KV cache $\gets \varnothing$
        \While{窗口数 $< L_{\max}$}
            \State $\mathbf{h} \gets M.\text{backbone}(x_{\text{last}},\ \text{KV cache})$
            \State $\text{logits} \gets M_{\mathrm{out}}(\mathbf{h})$
            \State $t \gets \text{sample}(\text{logits},\ T, k, p)$
            \State 将 $t$ 追加到 $x$
            \If{$t = \text{END}$} \State \textbf{break} \EndIf
        \EndWhile
        \State \Return $x$
    \end{algorithmic}
\end{algorithm}

当模型输出 END 时立即终止，如未输出 END 而达到最大窗口数则强制截断。

推理阶段还面临长度控制问题。若模型过早输出 END，序列可能只包含不完整面段或缺少边段；若迟迟不输出 END，则序列会达到最大长度后被截断，解析时也容易失败。本文把这两类情况都交由三态统计反映，而不是在采样阶段手写复杂规则强制补全。这样可以更真实地评估模型本身学到序列边界的程度。

\subsection{特殊符号与索引约束}

由于本文使用 $W = 1$ 的标准 next-token 头，原始 token 序列里的“面索引、坐标、特殊符号”三类 token 由模型按位置依次生成，面段、边段的分隔位置以及 SEP 的出现位置由训练数据中的高频统计模式自然学到。

\subsection{生成序列解析校验}

推理输出进入逆向重建前，会先经过轻量级解析防护。解析器检查 START、SEP、END 是否按预期出现，判断面段与边段长度能否分别被 49 和 14 整除，并验证面索引 token 是否落在 $[0,49]$ 范围内、坐标 token 是否落在量化坐标段。对于不满足基本格式的序列，重建程序不会强行补齐或猜测缺失字段，而是直接记录为 S1。这样处理虽然会降低尝试级有效率，但能避免把错误序列“修”成不可解释的实体，从而保持评测口径清晰。

本文没有在采样阶段加入强语法约束，而是在解析阶段做后验检查，主要出于两点考虑。第一，强约束解码需要显式维护当前位置属于面索引、坐标还是特殊符号，并在 softmax 后屏蔽大量非法 token，这会改变模型原始分布，使得生成有效率同时受解码规则影响。第二，本文希望先观察标准 next-token 模型在没有复杂外部规则时能够学到多少结构规律。hr2\_patch1 的 S1 为 0，说明在当前数据和采样配置下，模型已经能够较稳定地产生可解析序列，后续若面向更复杂数据集，再引入语法约束解码会更有针对性。

解析防护还为失败归因提供了一道清晰的边界。若样本在该阶段就被记为 S1，说明问题主要发生在序列语法或 token 类型层面，模型尚未稳定输出可解析的骨架，若通过解析但曲面拟合失败，问题转移到连续几何层面，更可能是采样点退化或拟合算法的容差不足，若已成功写出 STEP 但 BRepCheck 仍判无效，则问题进一步集中到拓扑缝合与内核合法性层面，最可能由相邻面边界不一致或面取向冲突造成。第~5 章的 S1 / S2 / S3 三态统计正是基于这种分层边界建立的，而本节描述的解析防护是这一分层得以稳定执行的工程基础——没有它，所有失败都会被混在同一个“重建失败”分类中，本文也就无法在结论中分别讨论哪一层的改进最为优先。

\section{序列到 B-rep 的逆向重建}

逆向重建管线由序列解析、NURBS 拟合、拓扑缝合与几何有效性检查组成，整个流程分为四步。

\begin{figure}[!htbp]
    \centering
    \includegraphics[width=0.82\textwidth]{插图/fig4-4_reconstruction_pipeline.pdf}
    \caption{逆向重建与评测管线示意图。生成序列首先按 START / SEP / END 被解析为面段与边段，并通过反量化恢复为三维采样点，随后分别拟合 NURBS 曲面和曲线，依据面索引关系完成拓扑缝合、实体构建、STEP / STL 导出与 BRepCheck 校验。图中下方同时给出 S1 重建失败、S2 产出但无效和 S3 产出且有效三类统计状态。}
    \label{fig:reconstruction-pipeline}
\end{figure}

\subsection{序列解析与坐标反量化}

解析器从 START 起按 SEP 切分面段与边段，再按 $B_f = 49$ 与 $B_e = 14$ 切分每个面与每条边。坐标 token 由反量化函数

\begin{equation}
    x = \frac{2 (t - \text{offset})}{N_q} - 1
    \label{eq:dequantization}
\end{equation}

恢复为 $[-1, 1]$ 内的浮点坐标。面索引段直接读取 token 数值作为面索引，特殊符号段触发段落切换或终止。

\subsection{NURBS 几何拟合}

对每个面，将 16 个反量化点送入 OpenCASCADE 的 B 样条曲面拟合接口，阶数取 $p = q = 3$，节点向量自动均匀分布，得到 NURBS 曲面；对每条边，将 4 个反量化点送入 B 样条曲线拟合接口，阶数取 3，得到 NURBS 曲线。拟合失败（如点共线、奇异参数）归入三态记录中的 S1（未产出）。

曲面和曲线拟合是从离散 token 回到 CAD 几何的关键环节。16 个面采样点并不等价于原始 STEP 中完整的 NURBS 控制网格，因此拟合结果更接近“由采样点重建出的近似面片”，而不是原始曲面的精确恢复。对于平面、圆柱面和规则过渡面，这种近似通常足以通过后续缝合，对于高曲率自由曲面或含小特征的面片，采样点数量可能不足以约束曲面形状，拟合误差会进一步传递到边界曲线和 BRepCheck。本文在讨论结果时将其视为整体重建管线误差的一部分，而不把神经网络输出与原始 CAD 参数完全等同。

边曲线拟合同样会影响拓扑闭合。若一条边的 4 个采样点与相邻两个面的边界位置不一致，Sewing 阶段就可能无法在容差范围内把面片缝合起来；即便后续 STEP writer 仍将 Shape 写出，BRepCheck 也会因边曲线不在相邻面上或开放边而把它归入 S2。由此可见，边段中两个邻接面索引和 12 个坐标 token 的正确性都很重要：前者决定边连接哪些面（拓扑信息），后者决定连接的位置与形状（几何信息）。本文点级表示之所以把这两类信息放在同一边段中，目的就是让模型在生成边时同步考虑拓扑邻接与几何走向。若先生成所有坐标再生成索引，模型很容易出现“几何上合理但索引错乱”或“索引正确但几何离面太远”两种典型失败模式；把两者交织在一起，可以使两类信息在注意力层面相互校准。

\subsection{拓扑缝合与实体构建}

\begin{sloppypar}
重建阶段先由拟合曲面创建面片，再对面集合执行缝合生成 Shell，随后使用形状修复工具处理部分容差不闭合的边，最后封闭为 Solid。实现中对应 OpenCASCADE 的建面、缝合、ShapeFix 与实体构造接口。Solid 经 OCC 的 STEP writer 导出 STEP 文件，并同时产出一份三角化的 STL 文件便于可视化。
\end{sloppypar}

该阶段是生成式 CAD 与一般三维点云生成区别最明显的部分。点云或网格生成只需产出一组点或三角面即可展示，而 B-rep 重建必须让曲面、曲线、边界和体拓扑在内核中形成一致结构。Sewing 操作尝试把相邻面片按容差缝合，ShapeFix 则对部分可修复的容差问题进行调整，但这些工具不能把任意错误序列变成合法实体。如果生成的邻接面索引混乱、边曲线远离相邻面，或面片本身退化严重，最终仍会在构造 Solid 或 BRepCheck 阶段失败。

因此，逆向重建既是输出转换模块，也是对前面生成质量的严格过滤器。本文没有把 ShapeFix 的自动修复能力视为模型性能的一部分来夸大，而是将最终 OCC 判定结果作为统一口径：能写出 STEP 且通过 BRepCheck 才进入 S3，写出但无效进入 S2，无法写出进入 S1。这样的统计方式可以清楚地区分“模型生成了结构上较完整的序列”和“后处理勉强写出了文件但仍不合法”两种情况。

\subsection{几何有效性校验}

最终调用 OpenCASCADE 的 BRepCheck 有效性检查作为判据，结果写入评测 JSON。有效性检查提供两个变体：一个从 STEP 路径读入并校验，另一个直接接收内存中的 Shape 实例进行校验。判据语义与第~2 章 §2.5 一致：内核校验通过即认定该样本“有效”。

\subsection{三态结果记录}

评测脚本输出五个标量：序列总数、重建成功、重建失败、有效 BREP、无效 BREP，按第~2 章\autoref{tab:three-state-def} 的三态口径分别对应 S1（重建失败）、S2（无效 BREP）、S3（有效 BREP）三个统计槽位。该输出格式直接支撑第~5 章主表口径。

五标量输出虽然简单，但对论文复核非常关键。只报告有效率百分比容易掩盖分母差异：例如，1000 次采样中 500 次无法写出 STEP、剩余 500 次全有效，与 1000 次采样全部写出且 500 次有效，在 of saved 口径下差异很大。本文同时记录采样次数、已写出 STEP 数和 BRepCheck Valid 数，可以让读者判断有效率来源。第~5 章对早期 hr2 基线和 hr2\_patch1 的比较也正依赖这一点：两者不仅 of saved 有效率不同，S1 重建失败比例也明显不同。
